% =============================================================================
% Introduction
% =============================================================================
\section{Introduction}\label{sec:introduction}

In 1966, Kac asked the now-classic question ``Can one hear the shape of a
drum?''~\cite{Kac1966}.  The question is spectral: does the set of resonant
frequencies determine the geometry uniquely?  Gordon, Webb, and
Wolpert~\cite{Gordon1992} later showed that, for planar drums, the answer is
``no'': distinct domains can share the same Laplacian spectrum.  That result
did not close the subject.  It made it more interesting.

Here we ask an engineering analogue.  Given the measured resonant frequencies
of a fluid-filled viscoelastic shell, can one recover its geometry and
stiffness?  More pointedly: can one hear the shape of an organ?  The question
is motivated by non-invasive tissue characterisation, vibroacoustic
elastography, and related inverse problems in biological and agricultural
materials~\cite{Sarvazyan1998,Doyley2012}.  If the answer is favourable, a
small set of resonant frequencies could carry useful information about organ
size, flattening, and effective stiffness.  The same logic applies, with only a
slightly different dinner plan, to fruit quality assessment.

The forward model considered here is the low-frequency flexural model developed
in our companion study~\cite{BrowntoneCompanion}.  We restrict attention to the
flexural family with mode number $n \geq 2$, rather than the breathing mode
$n=0$, which lies near \SI{2490}{Hz} and is governed primarily by fluid
compressibility.  The inverse problem is therefore posed on the measured
flexural frequencies
$\{f_2,f_3,\ldots,f_N\}$, with the aim of recovering the semi-major axis $a$,
the semi-minor axis $c$, and the Young's modulus $E$.

A central point must be stated plainly.  The implemented forward model is an
\emph{equivalent-sphere} model.  It uses the equivalent radius
\begin{equation}
  R_\mathrm{eq} = (a^2 c)^{1/3},
\end{equation}
and does not contain explicit spheroidal curvature terms or explicit
aspect-ratio corrections.  We therefore do \emph{not} present a full oblate
shell theory.  Instead, we examine the practical identifiability of the
equivalent-sphere model when parameterised by nominally oblate geometries
($a>c$).  Any geometric sensitivity enters only through
$R_\mathrm{eq}$ and through the mode-dependent balance of bending, membrane,
prestress, and added-mass contributions.

Our contribution is a practical identifiability analysis based on the Jacobian
of the frequency--parameter map.  We use the dimensionless scaled Jacobian,
its singular values, its condition number, and the associated Fisher
information matrix to assess whether small perturbations in the measured
frequencies can be mapped reliably back to $(a,c,E)$.  This is an analysis of
\emph{practical} identifiability in the sense of Raue et al.~\cite{Raue2009}
and Wieland et al.~\cite{Wieland2021}, not a global uniqueness theorem.

The main numerical findings are as follows.
\begin{itemize}
  \item At the spherical operating point, the scaled Jacobian is effectively
  rank deficient, with condition number
  $\kappa = 1.37 \times 10^{10}$.
  \item At the canonical oblate operating point, the same analysis gives
  $\kappa = 69.4$, an improvement by a factor of
  $2.0 \times 10^{8}$.
  \item A Newton-type inverse solver recovers $(a,c,E)$ to relative errors
  below $10^{-12}$ in round-trip tests when initialised within the local basin
  of attraction.
  \item At $1\%$ frequency noise, the Cram\'er--Rao lower bounds correspond to
  uncertainties of approximately $\pm 22.5\%$ in $a$, $\pm 6.0\%$ in $c$, and
  $\pm 30.8\%$ in $E$.
  \item Three-mode inversion is locally feasible but exhibits local minima,
  whereas five-mode inversion is robust across the tested initial guesses.
\end{itemize}

These results suggest that one can, in a practical sense, hear the shape of an
oblate organ from a modest set of flexural resonances, whereas the same task is
poorly posed at the spherical limit.  The distinction matters because many
inverse vibroacoustic schemes succeed or fail before the optimiser has even had
the chance to misbehave: the answer is already encoded in the conditioning of
the Jacobian.

The paper is organised as follows.  Section~\ref{sec:theory} summarises the
equivalent-sphere forward model and the inverse formulation.  Section
\ref{sec:methodology} describes the numerical Jacobian, the Newton-type
inversion procedure, and the parameter-space conditioning study.  Section
\ref{sec:results} presents the sphere--oblate contrast, singular-value
structure, Cram\'er--Rao bounds, and a watermelon cross-application.  Section
\ref{sec:discussion} relates the results back to Kac's question and to
vibroacoustic elastography, and Section~\ref{sec:conclusion} summarises the
practical recommendations.
